\section{Practical Appendices: Models, Templates, Checklists and Decision Exercises}

\subsection{Appendix A: The 1-Page Entrepreneurial Finance Dashboard}

Copy this structure into a weekly management note. Use actual numbers and mark unavailable fields honestly.

\begin{table*}[t]
\centering
\small
\caption{Appendix A: The 1-Page Entrepreneurial Finance Dashboard}
\label{tab:app_a_dashboard}
\begin{tabularx}{\textwidth}{p{2.4cm} p{4.6cm} p{5.0cm} X}
\toprule
\textbf{Domain} & \textbf{Metric or Field} & \textbf{Why It Matters} & \textbf{Review Frequency} \\
\midrule
Cash & Unrestricted bank cash reconciled to ledger & Actual liquidity & Weekly \\
Promises & Cost to fulfill prepaid and contracted obligations & Prevent spending service cash & Weekly \\
Survival & Minimum forecast cash in next 13 weeks & Identifies timing failure & Weekly \\
Sales & Cash collected versus revenue recognized & Explains income/cash difference & Monthly \\
Delivery & Contribution after all direct delivery costs & Establishes unit economics & Monthly \\
Founder subsidy & Founder hours and replacement value & Tests reproducibility & Monthly \\
Working capital & Receivable aging, WIP, payables, refunds & Measures trapped/released capital & Monthly \\
Capital & Capex approved, committed and actually spent & Prevents budget creep & Monthly \\
Debt & Debt service, interest reset, covenant headroom & Exposes fixed claims & Monthly \\
Ownership & Fully diluted table and outstanding convertibles & Keeps dilution visible & Per event \\
Reinvestment & Estimated incremental return by project & Directs scarce funds & Quarterly \\
Governance & Related-party exposure and approval status & Protects fiduciary integrity & Quarterly \\
Outcomes & Verified customer/learner performance & Keeps finance tied to value & Quarterly \\
\bottomrule
\end{tabularx}
\end{table*}

\subsection{Appendix B: 3-Statement Reconciliation Worksheet}

Reconcile statutory accounts and managerial economic models through a unified double-entry bridge:

\begin{tcolorbox}[colback=BayesCanvas,colframe=BayesDark,arc=1mm,title=\textbf{Appendix B: 3-Statement Reconciliation Worksheet}]
\scriptsize\ttfamily
Opening bank cash\\
+ Customer receipts\\
+ Other operating cash receipts\\
- Employee / supplier / tax operating cash paid\\
= Net operating cash flow (direct cash view)\\
- Capital expenditure + Asset disposal receipts\\
= Cash after investing activities\\
+ New equity and borrowing proceeds\\
- Principal repayment and owner distributions\\
= Net change in cash\\
+ Opening cash\\
= Closing cash (reconcile to bank and restricted-cash classification)\\[1ex]
Recognized revenue\\
- Recognized operating expense\\
- Depreciation and amortization\\
= Operating profit (subject to entity's presentation)\\
- Net finance expense\\
- Tax expense\\
= Net profit\\[1ex]
Opening equity + Net profit + New paid-in equity\\
- Dividends / repurchases +/- Other valid equity movements\\
= Closing equity\\[1ex]
Assets = Liabilities + Equity \quad [at each reporting date]
\end{tcolorbox}

\subsection{Appendix C: Founder-Adjusted Cohort Economics Worksheet}

Track the true marginal economics of cohorts by including full economic costs and replacement founder labor:

\begin{tcolorbox}[colback=BayesCanvas,colframe=BayesDark,arc=1mm,title=\textbf{Appendix C: Founder-Adjusted Cohort Economics Worksheet}]
\scriptsize\ttfamily
Paid seats                           N = \_\_\_\_\_\_\_\\
Realized price per paid seat         P = \_\_\_\_\_\_\_\\
Refunds / discounts / bad debts      R = \_\_\_\_\_\_\_\\
Net earned revenue                   N*P - R\\
External faculty and reviewer fees   F = \_\_\_\_\_\_\_\\
Delivery tools, cloud, facilities    T = \_\_\_\_\_\_\_\\
Acquisition and admissions           A = \_\_\_\_\_\_\_\\
Learner support and assessment       S = \_\_\_\_\_\_\_\\
Other incremental project expenses   O = \_\_\_\_\_\_\_\\
Founder teaching/support/sales hours H = \_\_\_\_\_\_\_\\
Replacement/opportunity rate/hour    W = \_\_\_\_\_\_\_\\
Cash contribution                    N*P - R - F - T - A - S - O\\
Founder-adjusted contribution        Cash contribution - H*W\\
Independent verified successes       Q = \_\_\_\_\_\_\_\\
Fully loaded cost/success            (F+T+A+S+O+H*W)/Q
\end{tcolorbox}

\subsection{Appendix D: NPV and Discount-Rate Consistency Checklist}

Verify that investment appraisal assumptions maintain economic consistency:

\begin{tcolorbox}[colback=BayesCanvas,colframe=BayesDark,arc=1mm,title=\textbf{Appendix D: NPV and Discount-Rate Consistency Checklist}]
\scriptsize\ttfamily
{[ ]} Incremental cash flows: baseline counterfactual defined (what happens if we do not invest?)\\
{[ ]} Cash flows are after-tax, after working capital and after maintenance capex\\
{[ ]} Financing flows excluded from unlevered cash flows if discounted at WACC\\
{[ ]} Currency of cash flows matches currency of discount rate\\
{[ ]} Nominal cash flows discounted at nominal rate; real cash flows at real rate\\
{[ ]} Opportunity costs of existing company-owned space, servers or faculty hours included\\
{[ ]} Cannibalization of existing revenue deducted as an incremental cost\\
{[ ]} Salvage or terminal value tested for sensitivity; not an arbitrary multiple\\
{[ ]} Downside cash trough identified (lowest liquidity point during rollout)\\
{[ ]} Discrete failure states evaluated via decision tree rather than single arbitrary rate
\end{tcolorbox}

\subsection{Appendix E: ROIC, Reinvestment and Economic Profit Calculator}

Evaluate whether operations generate economic surplus above the relevant cost of capital:

\begin{tcolorbox}[colback=BayesCanvas,colframe=BayesDark,arc=1mm,title=\textbf{Appendix E: ROIC, Reinvestment and Economic Profit Calculator}]
\scriptsize\ttfamily
Operating revenue                       200 lakh\\
Normalized EBIT                         30 lakh\\
Illustrative normalized tax rate        20\%\\
NOPAT                                    24 lakh\\
Average operating invested capital      120 lakh\\
ROIC                                     20\%\\
Relevant WACC                            12\%\\
ROIC-WACC spread                         8 percentage points\\
Economic profit                          24 - (120*0.12) = 9.6 lakh\\[1ex]
Reinvestment rate (stable illustration)  40\% of NOPAT\\
Return on incremental invested capital  20\%\\
Simplified sustainable operating growth  0.40*0.20 = 8\%
\end{tcolorbox}

\subsection{Appendix F: A Cash-Flow Scenario and Reverse-Stress Template}

Stress test operational liquidity across adverse scenarios and identify the boundary of solvency:

\begin{table*}[t]
\centering
\small
\caption{Appendix F: Cash-Flow Scenario and Reverse-Stress Testing Matrix}
\label{tab:app_f_stress}
\begin{tabularx}{\textwidth}{p{4.8cm} r r X}
\toprule
\textbf{Driver} & \textbf{Base Case} & \textbf{Downside Case} & \textbf{Severe Stress Case} \\
\midrule
Paid units or contracts & [ ] & [ ] & [ ] \\
Realized price & [ ] & [ ] & [ ] \\
Incremental servicing cost & [ ] & [ ] & [ ] \\
Required fixed capacity step & [ ] & [ ] & [ ] \\
Collections delay, days & [ ] & [ ] & [ ] \\
Contractual refunds or credits & [ ] & [ ] & [ ] \\
Maintenance and growth capex & [ ] & [ ] & [ ] \\
Debt service due & [ ] & [ ] & [ ] \\
Minimum monthly cash balance & [ ] & [ ] & [ ] \\
Minimum covenant headroom & [ ] & [ ] & [ ] \\
\bottomrule
\end{tabularx}
\end{table*}

Reverse-stress question: \textbf{What combination of demand, price and collection delay makes the minimum cash balance negative or breaks an essential quality standard?} Work backward until the boundary becomes explicit. Define a management action that occurs \textbf{before} reaching it. Do not confuse a cash forecast's year-end total with its lowest point during the year.

\subsection{Appendix G: Investment Approval Memorandum}

Require a structured memo before approving capital expenditure or significant long-term commitments:

\begin{tcolorbox}[colback=BayesCanvas,colframe=BayesDark,arc=1mm,title=\textbf{Appendix G: Investment Approval Memorandum Template}]
\scriptsize\ttfamily
1. METADATA: Decision date / sponsor / independent reviewer\\
2. NEED: Customer, stakeholder or mission need\\
3. COUNTERFACTUAL: Current process and baseline alternative\\
4. CAPABILITY: What durable asset or capability will remain\\
5. CASH REQUIRED: Cash required now / later / at worst\\
6. CASH-FLOW SCHEDULE: Incremental fully loaded cash-flow schedule\\
7. REINVESTMENT: Working capital, maintenance and replacement requirements\\
8. RETURN: Expected NPV or other defensible return measure\\
9. HURDLE BASIS: Reason discount rate or risk treatment is appropriate\\
10. MISSION VALUE: Nonfinancial mission benefit, separately described\\
11. STATUTORY CHECKS: Critical regulatory, IP and accounting conditions\\
12. EVIDENCE: Main assumptions and evidence quality\\
13. SCENARIOS: Base case / downside / reverse stress\\
14. CAPACITY: Founder hours and management capacity consumed\\
15. REVERSIBILITY: Staged milestones and stop conditions\\
16. DISPLACED USES: Alternative projects displaced\\
17. GOVERNANCE: Conflicts of interest and approvals\\
18. VERDICT: Decision: fund / pilot / defer / redesign / reject\\
19. REASSESSMENT: Named reviewer, date of future post-audit
\end{tcolorbox}

The reviewer should be able to understand \textbf{why the project is attractive without speaking to its proponent}. A beautifully formatted memo is worthless if it cannot reconcile to a credible cash plan.

\subsection{Appendix H: Financing Comparison Matrix}

Compare funding instruments across dilution, cash claims, operational fit, and founder risk:

\begin{table*}[t]
\centering
\small
\caption{Appendix H: Strategic Financing Comparison Matrix}
\label{tab:app_h_financing_matrix}
\begin{tabularx}{\textwidth}{p{2.6cm} p{2.2cm} p{3.2cm} p{3.2cm} X}
\toprule
\textbf{Instrument} & \textbf{Immediate Dilution} & \textbf{Mandatory Cash Payment} & \textbf{Typical Suitable Need} & \textbf{Primary Founder Risk} \\
\midrule
Retained earnings & None & None external & Proven incremental reinvestment & Concentration, lost liquidity \\
Customer advance & None & Delivery/refund obligations & Paid validated service & Spending future-service cash \\
Trade credit & None & Due at supplier terms & Short operating cycle & Supplier disruption \\
Bank term debt & None directly & Interest + principal & Predictable productive asset & Default, security/guarantees \\
Revolving credit & None directly & Interest and cleanup terms & Seasonal working capital & Draw restrictions, renewal \\
Factoring & None directly & Fee / recourse depends & Collectible receivable & Disputes, costly recourse \\
Revenue-share & Usually none & Share of receipts & Observable revenue stream & Revenue skim during low margins \\
Angel / common equity & Yes & No contractual repayment & Uncertain early growth & Dilution, investor conflict \\
Venture preferred & Yes & Usually no routine debt service & High-upside scalable venture & Preferences, board rights, exits \\
Convertible note & At conversion & Debt rights may mature & Bridge to financing & Maturity + dilution uncertainty \\
US-style SAFE & At conversion & Generally no note coupon/maturity & US early venture financing & Accumulated conversion dilution \\
Growth equity / PE & Yes & Deal-dependent & Expansion or acquisition & Control, leverage, exit timing \\
Grant / donation & Usually none & Restrictions / clawback & Eligible mission or research & Compliance, restricted funds \\
Project finance & Sponsor equity + debt & Contracted project debt & Long-lived contracted asset & Completion, DSCR, recourse \\
\bottomrule
\end{tabularx}
\end{table*}

Instrument economics vary dramatically by contract, country and investor. This is a comparison map, not a legal statement about a particular loan or security.

\subsection{Appendix I: A Fully Diluted Cap-Table Working Example}

A simplified, purely illustrative sequence with no notes, SAFEs or preferences:

\begin{table*}[t]
\centering
\small
\caption{Appendix I: Fully Diluted Capitalization Table Sequential Evolution}
\label{tab:app_i_captable}
\begin{tabularx}{\textwidth}{p{3.2cm} r r r r X}
\toprule
\textbf{Holder} & \textbf{Founding} & \textbf{Round A (20\% new)} & \textbf{Round B (25\% new)} & \textbf{New 10\% Pool} & \textbf{Economic Notes} \\
\midrule
Founder & 100\% & 80\% & 60\% & 54\% & Diluted proportionally across rounds \\
Investor A & - & 20\% & 15\% & 13.5\% & Diluted by Round B and option pool \\
Investor B & - & - & 25\% & 22.5\% & Diluted by post-round option pool \\
New option pool & - & - & - & 10\% & Designated unallocated employee equity \\
\midrule
\textbf{Total} & \textbf{100\%} & \textbf{100\%} & \textbf{100\%} & \textbf{100\%} & Fully diluted pro forma capitalization \\
\bottomrule
\end{tabularx}
\end{table*}

To compare real financing, include promised grants, existing unallocated pool, future pool requirement, warrants, convertible instruments, anti-dilution and investors' participation rights. \textbf{The option-pool creation sequence above is not the same as a pre-money pool expansion.} The latter may shift a greater share of dilution toward existing owners.

\subsection{Appendix J: Liquidation Preference Waterfall Exercise}

Assume a single investor paid \textbf{\textnormal{₹}2 crore for 20\% preferred equity}, no debt, transaction costs or other claims, and a \textnormal{₹}5 crore distributable sale price:

\begin{table*}[t]
\centering
\small
\caption{Appendix J: Liquidation Preference Waterfall Comparison on a \textnormal{₹}5 Crore Distributable Exit}
\label{tab:app_j_waterfall}
\begin{tabularx}{\textwidth}{p{4.2cm} r r X}
\toprule
\textbf{Security Terms} & \textbf{Preferred Investor} & \textbf{All Common Holders} & \textbf{Distribution Mechanics \& Explanation} \\
\midrule
As-converted ordinary 20\% & ₹1.0 crore & ₹4.0 crore & Pure percentage pro rata sharing across proceeds \\
1$\times$ non-participating preference & ₹2.0 crore & ₹3.0 crore & Chooses ₹2.0 crore preference over ₹1.0 crore conversion \\
1$\times$ uncapped participating & ₹2.6 crore & ₹2.4 crore & ₹2.0 crore preference 1st, then 20\% of remaining ₹3.0 crore \\
\bottomrule
\end{tabularx}
\end{table*}

At a \textbf{\textnormal{₹}20 crore} distributable exit, the 1$\times$ \textbf{non-participating} investor would instead convert for \textbf{\textnormal{₹}4 crore}, exceeding the \textnormal{₹}2 crore preference. A participating instrument, if uncapped, could yield \textbf{\textnormal{₹}2 crore + 20\% of \textnormal{₹}18 crore = \textnormal{₹}5.6 crore} in this simplified structure. Real agreements may include caps, dividends, senior securities, conversion adjustments and other rights. Use executed documents to model the waterfall.

\textbf{A question worth asking:} How much of a plausible mediocre exit goes to common holders \textit{after} the full preference stack, fees, escrow, debts and taxes? It often matters more than the best-case outcome shown during fundraising.

\subsection{Appendix K: Debt-Service and Refinancing Stress Sheet}

Map debt obligations and covenant headroom across the debt lifecycle:

\begin{tcolorbox}[colback=BayesCanvas,colframe=BayesDark,arc=1mm,title=\textbf{Appendix K: Debt-Service and Refinancing Stress Sheet}]
\scriptsize\ttfamily
Loan opening principal:\\
Annual or quarterly rate and reset index:\\
Scheduled principal in each period:\\
Interest and fees in each period:\\
Cash available for debt service (document covenant definition):\\
DSCR by period = available cash / (interest + scheduled principal)\\
Minimum required covenant ratio:\\
Excess/shortfall by period:\\
Required cash reserve / restricted collateral:\\
Guarantees, security and assets at risk:\\
Maturity and balloon payment:\\
Committed refinancing sources, if any:\\
Stress case if refinancing is unavailable:
\end{tcolorbox}

\subsection{Appendix L: Forensic Accounting and Quality-of-Earnings Checklist}

Conduct forensic analysis by verifying accounting claims against underlying bank records and operational facts:

\begin{enumerate}
    \item \textbf{Cash reconciliation:} Does operating cash flow reconcile to reported EBITDA after normalizing working capital and taxes?
    \item \textbf{Revenue verification:} Are invoices matched to signed customer contracts, approved milestones, delivery logs, and bank deposits?
    \item \textbf{Receivables aging:} Are receivables older than 60 days provisioned appropriately? Is there an end-of-quarter billing spike?
    \item \textbf{Deferred revenue integrity:} Is prepaid customer cash ring-fenced to cover the full expected cost of completing delivery?
    \item \textbf{Capitalization scrutiny:} Are capitalized development costs compliant with applicable accounting criteria and supported by timesheets?
    \item \textbf{Related-party transactions:} Are transactions with founders, family members, or affiliated companies conducted at documented arm's-length terms?
    \item \textbf{Expense completeness:} Are founder compensation, contractor fees, software subscriptions, and maintenance expenses fully reflected?
    \item \textbf{Off-balance-sheet commitments:} Are facility leases, minimum volume purchase contracts, and debt guarantees fully documented?
    \item \textbf{Adjustment audit:} Are add-backs to EBITDA genuinely nonrecurring, or do they represent recurring operational costs of doing business?
    \item \textbf{Tax and statutory compliance:} Are GST/VAT filings, payroll taxes, withholding deductions, and annual statutory filings up to date?
\end{enumerate}

\subsection{Appendix M: Project-Finance Investment Preflight}

Before committing to asset-heavy, long-lived infrastructure, verify project-finance prerequisites:

\begin{enumerate}
    \item \textbf{Separable cash flows:} Does the project possess an observable, legally segregated cash flow independent of parent company subsidy?
    \item \textbf{Contractual offtake:} Are long-term customer purchase commitments, enrollment agreements, or service contracts enforceable?
    \item \textbf{Turnkey construction:} Is construction governed by fixed-price, date-certain EPC contracts with enforceable liquidated damages?
    \item \textbf{Permits and land title:} Are clear legal land title, environmental clearances, municipal approvals, and educational zoning secured?
    \item \textbf{Debt-service coverage:} Does the downside scenario maintain a minimum DSCR of at least 1.30$\times$ across all amortization periods?
    \item \textbf{Reserve accounts:} Are a 6-month debt service reserve account (DSRA) and major maintenance reserve account fully funded from sources?
    \item \textbf{Recourse boundaries:} Are sponsor completion guarantees, cost-overrun undertakings, and contingent liabilities capped in writing?
    \item \textbf{Refinancing exposure:} Does the debt amortize fully within the contracted asset life, avoiding high-risk balloon maturities?
    \item \textbf{Downside recovery:} What is the liquidation or alternative-use value of the physical asset if primary customer demand collapses?
\end{enumerate}

\subsection{Appendix N: Grant and Philanthropic Funding Dossier}

Maintain institutional stewardship over non-dilutive grant and donor funding:

\begin{tcolorbox}[colback=BayesCanvas,colframe=BayesDark,arc=1mm,title=\textbf{Appendix N: Grant and Philanthropic Funding Dossier}]
\scriptsize\ttfamily
Funding source and legal recipient eligibility:\\
Charitable / for-profit / research entity status:\\
Restricted purpose and permitted costs:\\
Matching-fund requirement and timing:\\
Payment in advance or reimbursement:\\
Reporting and audit duties:\\
IP, publication and data restrictions:\\
Staff and overhead recovery permitted:\\
Clawback, termination and reputational risks:\\
Related-party disclosures and approvals:\\
Long-run maintenance funding after grant ends:\\
Independent mission-outcome evidence:
\end{tcolorbox}

\subsection{Appendix O: Capital Provider Interview Guide}

Interview prospective investors and lenders to evaluate incentive alignment and behavioral conduct under stress:

\begin{enumerate}
    \item \textbf{Fund economics:} What is the fund's vintage, total fund size, deployment pace, and standard follow-on reserve allocation?
    \item \textbf{Decision process:} Who serves on the investment committee, what are the approval hurdles, and what diligence is required?
    \item \textbf{Target outcome distribution:} What return multiple and exit timeline does your partnership model require from this specific investment?
    \item \textbf{Adversity conduct:} How did your firm support portfolio companies during severe macro downturns or operational missteps?
    \item \textbf{Founder references:} May we speak with founders of companies in your portfolio that missed projections or required restructuring?
    \item \textbf{Governance style:} What reserved matters, veto rights, board seats, and reporting formats do your standard term sheets require?
    \item \textbf{Value add reality:} Beyond capital, what concrete customer introductions, executive recruiting, or technical support do you deliver?
    \item \textbf{Follow-on policy:} Under what explicit conditions do you participate in insider bridges or lead follow-on financing rounds?
\end{enumerate}

\subsection{Appendix P: A 90-Day Financial-Capability Buildout}

Construct institutional financial capability in 6 disciplined 15-day sprints:

\begin{itemize}
    \item \textbf{Days 1-15: Establish cash baseline.} Segregate company bank accounts; reconcile bank balance to general ledger weekly; build a 13-week rolling cash forecast; ring-fence customer prepayment reserves.
    \item \textbf{Days 16-30: Normalize delivery economics.} Cost every direct course, pilot, or client engagement; track external contractor fees, server usage, and admissions spend; calculate founder labor replacement hours.
    \item \textbf{Days 31-45: Standardize monthly close.} Formalize revenue recognition rules; separate earned revenue from deferred contract liabilities; establish receivables aging and payables schedule; bridge profit to cash.
    \item \textbf{Days 46-60: Implement capital approval charter.} Adopt formal investment approval thresholds; require 1-page investment memos for capex; institute hurdle rate and alternative-use reviews.
    \item \textbf{Days 61-75: Stress financing structure.} If external financing is plausible, model debt service, dilution, options, preference and governance rights; request professional advice on valid instruments. Do not accept investment merely to complete the exercise.
    \item \textbf{Days 76-90: Institutionalize review.} Hold the 1st quarterly allocation review, compare earlier forecasts with outcomes and retire at least 1 weak assumption. Publish a management note that communicates uncertainty and sets the next evidence gate.
\end{itemize}

The completion test is not a beautiful spreadsheet. It is whether the founder can explain \textbf{where cash came from, why it differs from profit, which promises remain, which investments earn their cost, and what could make the business insolvent}.

\subsection{Appendix Q: Concept Map: Which Method Belongs to Which Question?}

Navigate corporate finance by matching analytical tools to the specific operating question:

\begin{table*}[t]
\centering
\small
\caption{Appendix Q: Capital Allocation and Corporate Finance Concept Map}
\label{tab:app_q_concept_map}
\begin{tabularx}{\textwidth}{p{4.2cm} p{5.0cm} X}
\toprule
\textbf{Founder Question} & \textbf{Primary Analytical Tools} & \textbf{Common Misuse \& Pitfall} \\
\midrule
Can we pay next month's obligations? & Cash forecast, payment register, reserves & Using EBITDA as cash \\
Does this sale contribute? & Incremental contribution and founder-hours costing & Ignoring support and acquisition \\
Does this investment add economic value? & Project FCFF, NPV, scenarios & Optimistic terminal value \\
How capital-efficient is our business? & ROIC and capital turnover & ROE without leverage context \\
How quickly can earnings grow sustainably? & Reinvestment runway $\times$ incremental returns & Assuming reinvestment opportunities never end \\
How much is this firm worth? & DCF, comparable economics, rights and claims & Equating paper valuation with liquidity \\
Should we borrow? & Downside cash and DSCR, covenant headroom & Borrowing because coupon looks cheap \\
Should we raise equity? & Funding milestones, cap-table scenarios, investor fit & Valuation alone \\
Do preferred terms matter? & Exit waterfall, protective-rights matrix & Nominal ownership percentage \\
Should we buy a facility? & Total project uses, life-cycle NPV, utilization & Ignoring operational readiness and maintenance \\
Are earnings trustworthy? & Quality-of-earnings and source-document review & Treating ratios as fraud verdicts \\
How should the board approve capital? & Thresholds, independent review, post-audits & Founder enthusiasm as substitute for evidence \\
\bottomrule
\end{tabularx}
\end{table*}

\subsection{Appendix R: Essential Finance Glossary for Private-Company Founders}

\textbf{Accrual accounting:} recognition of economic activity by period rather than cash exchange alone. \textbf{Accounts receivable:} amounts owed by customers. \textbf{Accounts payable:} amounts owed to suppliers. \textbf{Deferred revenue / contract liability:} cash or consideration received before performance required for revenue recognition. \textbf{Working capital:} resources tied to operating assets and operating liabilities. \textbf{Contribution margin:} revenue less explicitly defined incremental costs. \textbf{EBITDA:} earnings before interest, taxes, depreciation and amortization; not free cash flow. \textbf{EBIT:} earnings before interest and taxes. \textbf{NOPAT:} after-tax operating profit before financing effects. \textbf{FCFF:} after-tax free cash flow available to all capital providers. \textbf{Capex:} capital expenditure for durable assets. \textbf{Maintenance capex:} spending needed to sustain existing operations.

\textbf{NPV:} present value of incremental cash flows less initial investment. \textbf{IRR:} discount rate that makes project NPV zero. \textbf{Cost of equity:} expected return required by residual owners. \textbf{Cost of debt:} market cost of borrowing for relevant risk and terms. \textbf{WACC:} financing-weighted hurdle for appropriately matched unlevered cash flows. \textbf{ROIC:} NOPAT divided by average invested capital. \textbf{Economic profit:} NOPAT less a charge for invested capital. \textbf{Capital turnover:} revenue relative to invested operating capital. \textbf{DSCR:} cash available for debt service relative to contracted debt service under defined terms. \textbf{Leverage:} effect of fixed operating or financing commitments on residual outcomes.

\textbf{Pre-money valuation:} equity valuation immediately before new primary investment under an agreed capitalization definition. \textbf{Post-money valuation:} after that investment; other capitalization changes complicate interpretation. \textbf{Fully diluted ownership:} ownership including specified potential share issuances/conversions. \textbf{Option pool:} equity capacity designated for employee or other option grants. \textbf{SAFE:} contractual right to future equity under the relevant jurisdiction-specific instrument; a US post-money SAFE differs from a convertible note. \textbf{Convertible note:} debt with possible equity conversion. \textbf{Liquidation preference:} preferential payment claim in a specified exit or liquidation. \textbf{Participating preferred:} preferred claim that may also participate in residual proceeds. \textbf{Anti-dilution:} specified adjustment to conversion or share economics upon certain new issuances. \textbf{Pro rata right:} right to buy later securities to maintain an interest, under terms. \textbf{Drag/tag:} rights affecting participation in a sale of equity.

\textbf{Primary sale:} securities issuance that provides capital to the company. \textbf{Secondary sale:} transfer of an existing security whose proceeds typically go to the seller. \textbf{Project finance:} financing principally serviced by project cash flows through a structured project entity. \textbf{Restricted grant:} external mission funding with permitted-use conditions. \textbf{Quality of earnings:} analysis of the sustainability and reliability of reported profit. \textbf{Forensic accounting:} investigation of accounting evidence, transactions and control failures. \textbf{Related-party transaction:} exchange involving parties with special relationships requiring appropriate treatment and safeguards.

\section{Research Foundations and Official Reference Notes}

The following sources informed this chapter. They are \textbf{not a claim} that every proposed worksheet or Bayes recommendation is directly validated by research. Textbook models often require assumptions that fail in very early private companies. Regulatory sources are jurisdiction- and date-dependent; always consult current official provisions and qualified advisers.

\begin{enumerate}
    \item \textbf{Aswath Damodaran, NYU Stern: Corporate Finance resources.} Investment, financing, payout, valuation and the objective function; excellent open exercises and lectures. \url{https://pages.stern.nyu.edu/adamodar/New_Home_Page/CFin/CF.htm}
    \item \textbf{Aswath Damodaran: Introduction to Valuation.} Enterprise versus equity value, FCFF/FCFE, cost of capital and risk consistency. \url{https://pages.stern.nyu.edu/~adamodar/New_Home_Page/background/valintro.htm}
    \item \textbf{Aswath Damodaran: Investment Valuation resources.} Returns on capital, valuation, real options and industry dataset approaches. \url{https://pages.stern.nyu.edu/~adamodar/New_Home_Page/Inv2ed.htm}
    \item \textbf{Richard A. Brealey, Stewart C. Myers, Franklin Allen: \textit{Principles of Corporate Finance}.} Major textbook foundation for present value, investment choice, financing and capital structure. Consult an appropriate edition.
    \item \textbf{Jonathan Berk and Peter DeMarzo: \textit{Corporate Finance}.} Discounting, risk, capital budgeting, security structure and valuation. Consult an appropriate edition.
    \item \textbf{Franco Modigliani and Merton Miller (1958), `The Cost of Capital, Corporation Finance and the Theory of Investment.'} Benchmark for thinking about capital-structure irrelevance under deliberately restrictive assumptions and why actual taxes, frictions and distress matter. \url{https://www.jstor.org/stable/1809766}
    \item \textbf{Stewart C. Myers and Nicholas S. Majluf (1984), `Corporate Financing and Investment Decisions When Firms Have Information That Investors Do Not Have.'} Information asymmetry, financing choices and underinvestment. \url{https://doi.org/10.1016/0304-405X(84)90023-0}
    \item \textbf{Michael C. Jensen (1986), `Agency Costs of Free Cash Flow, Corporate Finance, and Takeovers.'} Caution about managerial discretion, free cash flow and capital governance. \url{https://www.jstor.org/stable/1818789}
    \item \textbf{Tim Koller, Marc Goedhart, David Wessels: \textit{Valuation: Measuring and Managing the Value of Companies}.} Value drivers, ROIC, growth, investment analysis and economic profit. Consult current edition and its assumptions.
    \item \textbf{William H. Beaver: \textit{Financial Reporting: An Accounting Revolution}.} Accounting information, reporting decisions and economic interpretation; consult an appropriate edition.
    \item \textbf{Stephen H. Penman: \textit{Financial Statement Analysis and Security Valuation}.} Accounting-quality and profitability analysis, reformulated statements and valuation. Consult an appropriate edition.
    \item \textbf{IFRS Foundation: IFRS 15 Revenue from Contracts with Customers.} Principles for recognition of revenue upon transfer of promised goods or services. \url{https://www.ifrs.org/issued-standards/list-of-standards/ifrs-15-revenue-from-contracts-with-customers/}
    \item \textbf{IFRS Foundation: IAS 7 Statement of Cash Flows.} Operating, investing and financing cash-flow classification and reconciliation. \url{https://www.ifrs.org/issued-standards/list-of-standards/ias-7-statement-of-cash-flows.html}
    \item \textbf{IFRS Foundation: IFRS 16 Leases.} Lessee recognition of many right-of-use assets and lease liabilities. \url{https://www.ifrs.org/issued-standards/list-of-standards/ifrs-16-leases/}
    \item \textbf{IFRS Foundation: IFRS 18 Presentation and Disclosure in Financial Statements.} Replaces IAS 1 for annual periods beginning on/after 1 January 2027 unless earlier applied. \url{https://www.ifrs.org/issued-standards/list-of-standards/ifrs-18-presentation-and-disclosure-in-financial-statements/}
    \item \textbf{Y Combinator: The SAFE: open-standard documentation and primer.} US post-money SAFE conversion, caps, MFN, no note-like maturity, rights and documentation; \textbf{not an India-law substitute}. \url{https://www.ycombinator.com/safe}
    \item \textbf{Y Combinator: SAFE vs. convertible note vs. priced round.} Instrument distinctions and common dilution errors. \url{https://www.ycombinator.com/safe/safe-vs-convertible-note}
    \item \textbf{World Bank PPP Reference: Project Finance: Key Concepts.} SPVs, nonrecourse/limited-recourse debt, contractual financing and construction cash-flow risk. \url{https://ppp.worldbank.org/financing/project-finance-concepts}
    \item \textbf{World Bank PPP Reference: Finance Structures for PPP.} Corporate versus project finance, recourse and asset security. \url{https://ppp.worldbank.org/finance-structures-ppp}
    \item \textbf{Institutional Limited Partners Association: ILPA Principles.} Alignment, transparency and governance in private-equity fund relationships; primarily LP/GP governance rather than startup term sheets. \url{https://ilpa.org/industry-guidance/principles-best-practices/ilpa-principles/}
    \item \textbf{Association of Certified Fraud Examiners: Occupational Fraud 2026: A Report to the Nations.} Occupational fraud case research and internal-control themes; case samples should not be treated as population prevalence estimates. \url{https://www.acfe.com/rttn}
    \item \textbf{Government of India, Department for Promotion of Industry and Internal Trade: Consolidated FDI Policy.} Relevant starting point for foreign investment and recognized-startup convertible notes; check amended FEMA rules and later directions as applicable. \url{https://www.dpiit.gov.in/static/uploads/2025/07/3ab2ec2a3bdb91c69653b7c34618c14a.pdf}
    \item \textbf{Government of India: The Companies Act, 2013 (India Code).} Corporate governance, share issuance, related-party and loan/investment provisions; consult current amendments, rules and authoritative consolidated text. \url{https://www.indiacode.nic.in/bitstream/123456789/15198/1/the_companies_act%2C_2013_no._18_of_2013_date_29.08.2013.pdf}
    \item \textbf{Reserve Bank of India: Master Direction on External Commercial Borrowings, Trade Credits and Structured Obligations.} Cross-border borrowing is regulated and changes over time; consult the RBI's current official directions for the actual borrower and lender. \url{https://www.rbi.org.in/}
    \item \textbf{DPIIT / Startup India official resources.} Recognition criteria and startup policies relevant to eligibility and specific instruments; verify the most current rules. \url{https://www.startupindia.gov.in/}
    \item \textbf{William F. Sharpe (1964), `Capital Asset Prices: A Theory of Market Equilibrium under Conditions of Risk.'} CAPM's equilibrium assumptions and the difference between systematic risk and entrepreneurial total risk. \url{https://doi.org/10.1111/j.1540-6261.1964.tb02865.x}
    \item \textbf{S. David Young and Stephen F. O'Byrne: \textit{EVA and Value-Based Management}.} Economic profit and capital-charge approaches; use adjusted definitions with care.
    \item \textbf{Jeffry Timmons and Stephen Spinelli: \textit{New Venture Creation}.} Entrepreneurial financing, venture economics and staged development. Consult an appropriate edition.
    \item \textbf{Brad Feld and Jason Mendelson: \textit{Venture Deals}.} Negotiating startup financing, preferred-share terms and venture governance. Consult a current edition; legal terms vary by jurisdiction.
    \item \textbf{Michael J. Mauboussin: capital allocation and expectations research.} Useful for the difference between growth and value creation and for expectations-sensitive investment decisions; consult published material with context.
\end{enumerate}

\subsection*{Suggested Learning Path Through the References}

Start with Damodaran's corporate-finance overview and a worked 3-statement example; learn NPV and cost-of-capital consistency; then study ROIC and reinvestment through valuation texts. Read the YC documentation when considering actual \textbf{US} early-stage financing, and lawyer-reviewed jurisdiction-specific documents when considering Indian transactions. Use World Bank project-finance materials only when a project has the contracted cash flows and legal structure necessary for such financing. Use IFRS and the local accounting standard appropriate to the entity to understand reporting quality. Treat forensic accounting as a discipline of evidence and corroboration, not an accusation engine.

\section*{Closing Principle}

\begin{quote}
\itshape
\textbf{Capital allocation is the founder's repeated decision about which future deserves to be made real.} Financing is the corresponding decision about who pays, who bears loss, who receives upside, and who gains the authority to influence that future. The 1st requires economic judgment; the 2nd requires contractual and institutional judgment; both require truthful accounting and financial restraint. The enduring organization is not the entity that raised the most money or owned the most assets. It is the institution that repeatedly \textbf{turned capital into valuable, reproducible capability without losing the right to continue.}
\end{quote}
